\thispagestyle{plain}
\addcontentsline{toc}{chapter}{General Conclusion}
\section*{General Conclusion}
This internship set out to address the infrastructure and deployment 
limitations of the Mashroom platform, which previously relied on a 
single virtual machine, manual deployments, and no environment 
separation. The solution designed and implemented over the course 
of this project replaces that setup with a containerized, 
cloud-native architecture running on Amazon EKS, provisioned 
entirely through Terraform, and delivered through a fully automated 
CI/CD pipeline combining GitHub Actions and ArgoCD. The result is 
a platform that is resilient, reproducible, and requires no manual 
intervention from commit to deployment.

Beyond the technical deliverables, this project was an opportunity 
to apply DevOps principles in a real-world setting, navigating the 
constraints of a production environment, making architectural 
decisions under practical limitations, and delivering a solution 
that is maintainable by the team long after the internship ends. 
The infrastructure is fully documented as code, the pipelines are 
self-explanatory by design, and the GitOps approach ensures that 
the cluster state is always traceable back to a repository commit.

That said, several aspects remain open for future improvement. 
The cluster currently has no observability tool in place for 
monitoring and alerting, making it difficult to detect and 
diagnose issues in production. Setting up a dedicated monitoring 
solution represents a natural next step toward a more mature 
and production-ready platform.